\documentclass[10pt]{extarticle}
\usepackage[a4paper,margin=1.5cm]{geometry}
\usepackage[T1]{fontenc}
\usepackage{amsmath,amssymb}
\usepackage{graphicx}
\usepackage[export]{adjustbox}
\usepackage{multicol}
\usepackage{enumitem}
\usepackage{booktabs,array,tabularx}
\setlength{\columnsep}{0.7cm}
\setlength{\parindent}{0pt}
\setlength{\parskip}{3pt}
\pagestyle{empty}
\makeatletter
\newcommand{\qfigdef}[4]{\@namedef{qf@file@#1}{#2}\@namedef{qf@w@#1}{#3}\@namedef{qf@wide@#1}{#4}}
\newcommand{\qf@check}[1]{\@ifundefined{qf@file@#1}{\PackageError{score2}{Unknown figure id #1}{}}{}}
\newcommand{\qf@img}[4][]{\includegraphics[width=#2\dimexpr\@nameuse{qf@w@#4}pt\relax,max width=#3,max height=0.45\textheight,#1]{\@nameuse{qf@file@#4}}}
\newif\ifqf@full
\newif\ifqf@boxed
\newsavebox{\qf@box}
\newcommand{\qf@open}{\global\qf@fulltrue\par\global\qf@boxedtrue
  \setbox\qf@box\vbox\bgroup\hsize\textwidth\linewidth\textwidth}
\newcommand{\qf@emit}[1]{\par\egroup\global\qf@boxedfalse
  \ifdim\dimexpr\ht\qf@box+\dp\qf@box\relax<#1\textheight\noindent\usebox{\qf@box}\par\else\unvbox\qf@box\fi}
\newcommand{\qfullsplit}{\ifqf@boxed\qf@emit{0.85}\fi}
\newcommand{\qf@close}{\ifqf@boxed\qf@emit{0.45}\fi\global\qf@fullfalse}
\newcommand{\qfullbegin}{\end{multicols}\qf@open}
\newcommand{\qfullend}{\qf@close\begin{multicols}{2}}
\newcommand{\qfullstart}{\qf@open}
\newcommand{\qfulldone}{\qf@close}
\newcommand{\qfig}[2][1]{\qf@check{#2}%
  \ifqf@full\qf@blk{#1}{#2}\else
  \ifnum\@nameuse{qf@wide@#2}=1 \expandafter\@firstoftwo\else\expandafter\@secondoftwo\fi
  {\end{multicols}\begin{center}\qf@img{#1}{\textwidth}{#2}\end{center}\begin{multicols}{2}}%
  {\qf@blk{#1}{#2}}\fi}
\newcommand{\qf@blk}[2]{\par\smallskip{\centering\qf@img{#1}{\linewidth}{#2}\par}\smallskip}
\newcommand{\qfigi}[2][1]{\qf@check{#2}\qf@img[valign=c]{#1}{\linewidth}{#2}}
\makeatother
\newcommand{\papersec}[1]{\par\bigskip{\centering\Large\bfseries PAPER-#1\par}\medskip}
\newcommand{\partsec}[2]{\par\medskip\noindent{\large\bfseries PART-#1: #2}\par\smallskip}
\newcommand{\sect}[1]{\par\smallskip\noindent\textbf{#1}\par}
\newcommand{\qn}[1]{\par\smallskip\noindent\textbf{#1.}~}
\newcommand{\paraq}[2]{\par\smallskip\noindent\textbf{Paragraph for Questions #1 and #2}\par}
\newenvironment{qoptions}{\begin{itemize}[leftmargin=1.8em,itemsep=1pt,topsep=2pt,label={}]}{\end{itemize}}
\newcommand{\opt}[2]{\item[(#1)] #2}
% figure registry (generated by score2_tools.py finalize)
\qfigdef{p1_phy_s1_q01_1}{Score-II - Test 45 (26 Mar 2025)-img/p1_phy_s1_q01_1.png}{75.9}{0}
\qfigdef{p1_phy_s1_q02_1}{Score-II - Test 45 (26 Mar 2025)-img/p1_phy_s1_q02_1.png}{50.4}{0}
\qfigdef{p1_phy_s1_q04_1}{Score-II - Test 45 (26 Mar 2025)-img/p1_phy_s1_q04_1.png}{91.4}{0}
\qfigdef{p1_phy_s1_q05_1}{Score-II - Test 45 (26 Mar 2025)-img/p1_phy_s1_q05_1.png}{91.6}{0}
\qfigdef{p1_phy_s1_q07_1}{Score-II - Test 45 (26 Mar 2025)-img/p1_phy_s1_q07_1.png}{100.4}{0}
\qfigdef{p1_phy_s1_q08_1}{Score-II - Test 45 (26 Mar 2025)-img/p1_phy_s1_q08_1.png}{117.4}{0}
\qfigdef{p1_phy_s1_q10_1}{Score-II - Test 45 (26 Mar 2025)-img/p1_phy_s1_q10_1.png}{86.3}{0}
\qfigdef{p1_phy_s1_q17_1}{Score-II - Test 45 (26 Mar 2025)-img/p1_phy_s1_q17_1.png}{104.6}{0}
\qfigdef{p1_phy_s1_q20_1}{Score-II - Test 45 (26 Mar 2025)-img/p1_phy_s1_q20_1.png}{88.8}{0}
\qfigdef{p1_phy_s2_q01_1}{Score-II - Test 45 (26 Mar 2025)-img/p1_phy_s2_q01_1.png}{150.0}{0}
\qfigdef{p1_chem_s1_q01_1}{Score-II - Test 45 (26 Mar 2025)-img/p1_chem_s1_q01_1.png}{95.1}{0}
\qfigdef{p1_chem_s1_q03_1}{Score-II - Test 45 (26 Mar 2025)-img/p1_chem_s1_q03_1.png}{25.2}{0}
\qfigdef{p1_chem_s1_q05_1}{Score-II - Test 45 (26 Mar 2025)-img/p1_chem_s1_q05_1.png}{107.1}{0}
\qfigdef{p1_chem_s1_q05_2}{Score-II - Test 45 (26 Mar 2025)-img/p1_chem_s1_q05_2.png}{151.0}{0}
\qfigdef{p1_chem_s1_q10_1}{Score-II - Test 45 (26 Mar 2025)-img/p1_chem_s1_q10_1.png}{56.1}{0}
\qfigdef{p1_chem_s1_q10_2}{Score-II - Test 45 (26 Mar 2025)-img/p1_chem_s1_q10_2.png}{50.8}{0}
\qfigdef{p1_chem_s1_q10_3}{Score-II - Test 45 (26 Mar 2025)-img/p1_chem_s1_q10_3.png}{59.1}{0}
\qfigdef{p1_chem_s1_q10_4}{Score-II - Test 45 (26 Mar 2025)-img/p1_chem_s1_q10_4.png}{54.2}{0}
\qfigdef{p1_chem_s1_q12_1}{Score-II - Test 45 (26 Mar 2025)-img/p1_chem_s1_q12_1.png}{78.5}{0}
\qfigdef{p1_chem_s1_q12_2}{Score-II - Test 45 (26 Mar 2025)-img/p1_chem_s1_q12_2.png}{78.5}{0}
\qfigdef{p1_chem_s1_q12_3}{Score-II - Test 45 (26 Mar 2025)-img/p1_chem_s1_q12_3.png}{78.5}{0}
\qfigdef{p1_chem_s1_q12_4}{Score-II - Test 45 (26 Mar 2025)-img/p1_chem_s1_q12_4.png}{78.5}{0}
\qfigdef{p1_chem_s2_q02_1}{Score-II - Test 45 (26 Mar 2025)-img/p1_chem_s2_q02_1.png}{107.8}{0}
\qfigdef{p1_chem_s2_q04_1}{Score-II - Test 45 (26 Mar 2025)-img/p1_chem_s2_q04_1.png}{41.6}{0}
\makeatletter\@beginparpenalty=10000\makeatother
\begin{document}
\begin{center}{\Large\bfseries JEE (Main+Advanced) Score-II --- Test 45 (26 Mar 2025)}\end{center}
\begin{multicols}{2}
\raggedcolumns
\partsec{1}{PHYSICS}
\sect{Section-I}

\qn{1} For what value of $R_2$, power loss through $R_1$ will be maximum?
\qfig{p1_phy_s1_q01_1}
\begin{qoptions}
\opt{A}{$R_1=R_2$}
\opt{B}{$R_2=2R_1$}
\opt{C}{$R_2=0$}
\opt{D}{$R_2=4R_1$}
\end{qoptions}

\qn{2} A rod AB of length $L$ and mass $M$ is uniformly charged with a charge $Q$, and is suspended from end A as shown in figure. The rod can freely rotate about A in the plane of the figure. A uniform electric field $E$ is switched on in the horizontal direction due to which the rod gets turned by a maximum angle of 90$^\text{o}$. The magnitude of $E$ is equal to
\qfig{p1_phy_s1_q02_1}
\begin{qoptions}
\opt{A}{$\dfrac{Mg}Q$}
\opt{B}{$\dfrac{Mg}{2Q}$}
\opt{C}{$\dfrac{2Mg}Q$}
\opt{D}{$\dfrac{4Mg}Q$}
\end{qoptions}

\qn{3} \textbf{Assertion (A) :} When a longitudinal pressure wave is reflected at the open end of an organ pipe, the compression pressure wave pulse becomes rarefaction pressure wave pulse during the reflection.\\
\textbf{Reason (R) :} The phase of the wave changes by $\pi$ when reflected at the open end.
\begin{qoptions}
\opt{A}{Both A and R are true and R is correct explanation of A.}
\opt{B}{Both A and R are true but R is not correct explanation of A.}
\opt{C}{A is true but R is false.}
\opt{D}{A is false but R is true.}
\end{qoptions}

\qn{4} The capacitance of the system is $C$. If the key is closed, the total energy loss is equal to :
\qfig{p1_phy_s1_q04_1}
\begin{qoptions}
\opt{A}{$\dfrac{Q^2}C$}
\opt{B}{$\dfrac{2Q^2}C$}
\opt{C}{$\dfrac{9Q^2}{8C}$}
\opt{D}{None of these}
\end{qoptions}

\qn{5} In the circuit shown, a conducting wire CD is moved with constant speed $v$ towards left. The complete circuit is placed in a uniform magnetic field $\vec B$, perpendicular to the plane of the circuit in inward direction. Current in CEFD is :-
\qfig{p1_phy_s1_q05_1}
\begin{qoptions}
\opt{A}{Clock wise}
\opt{B}{Anti-clockwise}
\opt{C}{Alternating}
\opt{D}{Zero}
\end{qoptions}

\qn{6} The magnetic field of a plane electromagnetic wave is given by $\vec B=3\times10^{-8}\sin[200\pi\,(y+ct)]\hat i$ T, where $c=3\times10^8$ m/s. The corresponding electric field is :-
\begin{qoptions}
\opt{A}{$\vec E=9\sin[200\pi\,(y+ct)]\hat k$ V/m}
\opt{B}{$\vec E=-10^{-6}\sin[200\pi\,(y+ct)]\hat j$ V/m}
\opt{C}{$\vec E=3\times10^{-8}\sin[200\pi\,(y+ct)]\hat k$ V /m}
\opt{D}{$\vec E=-9\sin[200\pi\,(y+ct)]\hat k$ V/m}
\end{qoptions}

\qn{7} In the circuit with ideal diodes as shown, current (in A) through battery is :-
\qfig{p1_phy_s1_q07_1}
\begin{qoptions}
\opt{A}{3}
\opt{B}{5}
\opt{C}{6}
\opt{D}{4}
\end{qoptions}

\qn{8} The centre of mass of disc of radius $\dfrac8{\sqrt5}$ m is moving with a velocity of 4 m/s on a horizontal surface. Its angular velocity is $\sqrt5$ rad/sec. The radius of curvature of point P as shown in figure is :-
\qfig{p1_phy_s1_q08_1}
\begin{qoptions}
\opt{A}{2m}
\opt{B}{2.5m}
\opt{C}{4m}
\opt{D}{5m}
\end{qoptions}

\qn{9} In elliptical orbit of a planet, as the planet moves from aphelion position to perihelion position, match the following :-
\par\nobreak
\begin{center}\small
\begin{tabular}{|c|p{3.0cm}|c|p{1.6cm}|}\hline
 & \multicolumn{1}{c|}{\textbf{Column-I}} & & \multicolumn{1}{c|}{\textbf{Column-II}}\\\hline
(P) & Speed of planet & (1) & Remains same\\\hline
(Q) & Distance of planet from center of sun & (2) & Decreases\\\hline
(R) & Potential energy & (3) & Increases\\\hline
(S) & Angular momentum of planet about the centre of the sun & (4) & Can not say\\\hline
\end{tabular}
\end{center}
\begin{qoptions}
\opt{A}{P $\to$ 3; Q $\to$ 4; R $\to$ 2; S $\to$ 1}
\opt{B}{P $\to$ 3; Q $\to$ 4; R $\to$ 2; S $\to$ 3}
\opt{C}{P $\to$ 3; Q $\to$ 1; R $\to$ 3; S $\to$ 2}
\opt{D}{P $\to$ 3; Q $\to$ 2; R $\to$ 2; S $\to$ 1}
\end{qoptions}

\qn{10} A point object is placed at a distance of 60 cm from a convex lens of focal length 30 cm. If a plane mirror is placed perpendicular to the principal axis of the lens at a distance of 40 cm from it, the final image would be formed at a distance of
\qfig{p1_phy_s1_q10_1}
\begin{qoptions}
\opt{A}{20 cm from the lens, it would be a real image.}
\opt{B}{30 cm from the lens, it would be a real image.}
\opt{C}{30 cm from the plane mirror, it would be a virtual image.}
\opt{D}{20 cm from the plane mirror, it would be a virtual image.}
\end{qoptions}

\qn{11} A copper wire and a steel wire of the same diameter and length are connected end to end. A force is applied which stretches their combined length by 1 cm. Then the wires have:
\begin{qoptions}
\opt{A}{the same stress and strain}
\opt{B}{the same stress but different strain}
\opt{C}{the same strain but different stress}
\opt{D}{different stress and strain}
\end{qoptions}

\qn{12} 20 gm of ice at $-40^\circ$C is mixed with 4 gm of steam at 100$^\circ$C. The final temperature of mixture is approximately (Latent heat of fusion is 80 cal/gm, Latent heat of vaporisation is 540 cal/gm, Specific heat of water is 1 cal/gm-$^\circ$C, specific heat of ice is 0.5 cal/gm$^\circ$C)
\begin{qoptions}
\opt{A}{30$^\circ$C}
\opt{B}{23$^\circ$C}
\opt{C}{28$^\circ$C}
\opt{D}{35$^\circ$C}
\end{qoptions}

\qn{13} A uniform string of length `$\ell$' is suspended from a rigid support. A short wave pulse is introduced at its lower end. It starts moving up the string. Choose the correct option.
\begin{qoptions}
\opt{A}{Acceleration of pulse increases}
\opt{B}{Acceleration of pulse decreases}
\opt{C}{Acceleration of pulse remains constant}
\opt{D}{None of the above}
\end{qoptions}

\qn{14} A plane polarised electromagnetic wave is incident on a medium of refractive index $\sqrt3$ from air. There is no partially reflected wave. So the angle of refraction of electromagnetic wave is :-
\begin{qoptions}
\opt{A}{45$^\text{o}$}
\opt{B}{60$^\text{o}$}
\opt{C}{30$^\text{o}$}
\opt{D}{90$^\text{o}$}
\end{qoptions}

\qn{15} \textbf{Statement-I :} If Gaussian surface does not enclose any charge, then electric field at any point on the Gaussian surface must be zero.\\
\textbf{Statement-II :} No net charge enclosed by Gaussian surface means net electric flux passing through the Gaussian surface is zero.
\begin{qoptions}
\opt{A}{Both Statement I and II are correct.}
\opt{B}{Statement I is correct but statement II is incorrect.}
\opt{C}{Statement I is incorrect but statement II is correct.}
\opt{D}{Both Statement I and II are incorrect.}
\end{qoptions}

\qn{16} If two SHM's whose equations are $x_1=3\sin\omega t$ and $x_2=-4\cos\omega t$ along the same line are superimposed, the equation of resultant SHM is :-
\begin{qoptions}
\opt{A}{$\sin(\omega t)$}
\opt{B}{$5\sin(\omega t-53^\text{o})$}
\opt{C}{$5\sin(\omega t+53^\text{o})$}
\opt{D}{$5\sin(\omega t-37^\text{o})$}
\end{qoptions}

\qn{17} In the given AC circuit $V_1$ \& $V_2$ are ideal voltmeter. The ratio of reading of $V_1$ to $V_2$.
\qfig{p1_phy_s1_q17_1}
\begin{qoptions}
\opt{A}{$\dfrac12$}
\opt{B}{$\dfrac32$}
\opt{C}{$\dfrac52$}
\opt{D}{$\dfrac34$}
\end{qoptions}

\qn{18} In a young's double slit experiment, light of wavelength 600 nm is used. P is the nearest point from the central maxima where intensity is $\dfrac34$th of maximum intensity. What percentage of the maximum intensity is observed at P if light is replaced by the light of wavelength 400 nm?
\begin{qoptions}
\opt{A}{50\%}
\opt{B}{75\%}
\opt{C}{25\%}
\opt{D}{100\%}
\end{qoptions}

\qn{19} A paramagnetic material has $10^{28}$ atoms/m$^3$. Its magnetic susceptibility at temperature 350 K is $2.8\times10^{-4}$. Its magnetic susceptibility at 300 K is :-
\begin{qoptions}
\opt{A}{$3.672\times10^{-4}$}
\opt{B}{$3.726\times10^{-4}$}
\opt{C}{$3.267\times10^{-4}$}
\opt{D}{$2.672\times10^{-4}$}
\end{qoptions}

\qn{20} The variation of the speed with time, of a ball falling in air is shown in below.
\qfig{p1_phy_s1_q20_1}
During the time from 0 to $T$, the ball gains some kinetic energy and loses gravitational potential energy $\Delta E_P$. Which of the following statements must be correct?
\begin{qoptions}
\opt{A}{$\Delta E_P$ is equal to gain in kinetic energy}
\opt{B}{$\Delta E_P$ is equal to the work done against air resistance}
\opt{C}{$\Delta E_P$ is greater than the gain in kinetic energy}
\opt{D}{$\Delta E_P$ is less than the work done against air resistance}
\end{qoptions}

\sect{Section-II}

\qn{1} The blocks are kept on smooth horizontal surface. If $N_1$ and $N_2$ are contact force between blocks for figure (i) \& (ii) respectively, then $\dfrac{2N_1}{N_2}$ is :-
\qfig{p1_phy_s2_q01_1}

\qn{2} A liquid of density 750 kg m$^{-3}$ flows smoothly through a horizontal pipe that tapers in cross-sectional area from $A_1=1.2\times10^{-2}$ m$^2$ to $A_2=\dfrac{A_1}2$. The pressure difference between the wide and narrow section of pipe is 4500 Pa. The rate of flow of liquid is $n\times10^{-3}$ m$^3$ s$^{-1}$, find $n$.

\qn{3} Radiation of photon energy 5eV is falling on a photo sensitive surface of work function 3eV. Consider an e$^-$ which makes two collisions before coming out of the surface and there is a loss of 20\% of energy in each collision. If K.E. of photoelectron coming out of surface is $x$eV, then $100x$ is :-

\qn{4} An ideal monoatomic gas undergoes a process in which its internal energy $U$ and density $\rho$ are related as $U\rho=$ constant. The ratio of change in internal energy and the work done by gas is $0.5x$. Find the value of $x$.

\qn{5} The dimensions of the rectangular block measured with a vernier caliper having least count of 0.1 mm is 5 mm $\times$ 10 mm $\times$ 5 mm. The maximum percentage error in measurement of volume of block is :-

\partsec{2}{CHEMISTRY}
\sect{Section-I}

\qn{1} The number of optically active isomers formed in the following reaction is
\qfig{p1_chem_s1_q01_1}
\begin{qoptions}
\opt{A}{2}
\opt{B}{4}
\opt{C}{6}
\opt{D}{7}
\end{qoptions}

\qn{2} Among the following, the correct statement for thionyl tetrafluoride is
\begin{qoptions}
\opt{A}{the geometry of thionyl tetrafluoride is trigonal bipyramidal having the sulphur-oxygen bond on the trigonal plane.}
\opt{B}{the geometry of thionyl tetrafluoride is trigonal bipyramidal having the sulphur-oxygen bond perpendicular to the trigonal plane.}
\opt{C}{the geometry of thionyl tetrafluoride is square pyramidal having the sulphur-oxygen bond on the square plane.}
\opt{D}{the geometry of thionyl tetrafluoride is square pyramidal having the sulphur-oxygen bond perpendicular to the square plane.}
\end{qoptions}

\qn{3} \textbf{Statement I :-} Compound \qfigi{p1_chem_s1_q03_1} is antiaromatic in nature.\\
\textbf{Statement II :-} Compounds having (4n) $\pi$-electrons and planer with a cyclic conjugated $\pi$-system is anti aromatic.
\begin{qoptions}
\opt{A}{Statement I is false but Statement II is true}
\opt{B}{Both Statement I and Statement II are true}
\opt{C}{Both Statement I and Statement II are false}
\opt{D}{Statement I is true but Statement II is false}
\end{qoptions}

\qn{4} A pink coloured aqueous solution of Co(NO$_3$)$_2$ turns blue on addition of HCl gradually. This colour change happens due to the formation of
\begin{qoptions}
\opt{A}{$\left[\mathrm{CoCl_4}\right]^{2-}$}
\opt{B}{$\left[\mathrm{CoCl_6}\right]^{4-}$}
\opt{C}{$\left[\mathrm{Co(H_2O)_4Cl_2}\right]$}
\opt{D}{$\left[\mathrm{Co(H_2O)_2Cl_4}\right]^{2-}$}
\end{qoptions}

\qn{5} Consider the following reaction sequence
\qfig{p1_chem_s1_q05_1}
\qfig{p1_chem_s1_q05_2}
How many sp$^3$ carbons are present in product (D), (assuming product (C) is benzene)
\begin{qoptions}
\opt{A}{7}
\opt{B}{8}
\opt{C}{6}
\opt{D}{9}
\end{qoptions}

\qn{6} Three different electrolytic cells containing AgNO$_3$, CuSO$_4$, AuCl$_3$, respectively, are electrolysed. Find the simplest ratio of mass of metal deposited if the cells are connected in parallel and the ratio of resistances is 3 : 2 : 1 respectively.\\
given :- Mw(Ag) = 108,\\
Mw(Cu) = 64 \& Mw(Au) = 197
\begin{qoptions}
\opt{A}{208 : 48 ; 197}
\opt{B}{108 : 48 : 197}
\opt{C}{108 : 64 : 197}
\opt{D}{208 : 64 : 197}
\end{qoptions}

\qn{7} Which of the following compound gives only one product on ozonolysis ?
\begin{qoptions}
\opt{A}{Ethene}
\opt{B}{Propene}
\opt{C}{1-Butene}
\opt{D}{2-Methyl Propene}
\end{qoptions}

\qn{8} Which of the following compound does not undergo a nucleophilic substituion reation?
\begin{qoptions}
\opt{A}{1-Chlorobutene}
\opt{B}{Chlorobenzene}
\opt{C}{Benzyl chloride}
\opt{D}{2-Chloropropane}
\end{qoptions}

\qn{9} Match the column and choose the correct option
\par\nobreak
\begin{center}\footnotesize
\begin{tabular}{|c|p{2.6cm}|c|p{2.5cm}|}\hline
\multicolumn{2}{|c|}{\shortstack{\textbf{Column-I}\\\textbf{(Properties)}}} & \multicolumn{2}{c|}{\shortstack{\textbf{Column-II}\\\textbf{(Order)}}}\\\hline
(A) & Electronegativity & (1) & B $<$ C $<$ N $<$ O\\\hline
(B) & Cationic size & (2) & Li $>$ Mg $>$ Be\\\hline
(C) & Metallic Character & (3) & K $>$ Mg $>$ Al\\\hline
(D) & Electron affinity & (4) & Cl $>$ F $>$ Br $>$ I\\\hline
\end{tabular}
\end{center}
\begin{qoptions}
\opt{A}{A $-$ 1, B $-$ 2, C $-$ 3, D $-$ 4}
\opt{B}{A $-$ 4, B $-$ 3, C $-$ 2, D $-$ 1}
\opt{C}{A $-$ 2, B $-$ 3, C $-$ 4, D $-$ 1}
\opt{D}{A $-$ 3, B $-$ 2, C $-$ 4, D $-$ 1}
\end{qoptions}

\qn{10} Which of the following compounds has the highest number of chiral centres?
\begin{qoptions}
\opt{A}{\qfigi{p1_chem_s1_q10_1}}
\opt{B}{\qfigi{p1_chem_s1_q10_2}}
\opt{C}{\qfigi{p1_chem_s1_q10_3}}
\opt{D}{\qfigi{p1_chem_s1_q10_4}}
\end{qoptions}

\qn{11} The magnetic moments in BM and geometry of complex $\left[\mathrm{Ni(CN)_4}\right]^{4-}$
\begin{qoptions}
\opt{A}{0.0 , TD}
\opt{B}{3.87, TBP.}
\opt{C}{1.73, SP}
\opt{D}{3.87, SP.}
\end{qoptions}

\qn{12} Which of the following plot represent positive deviation in a binary solution of liquids `A' and `B'.
\begin{qoptions}
\opt{A}{\qfigi{p1_chem_s1_q12_1}}
\opt{B}{\qfigi{p1_chem_s1_q12_2}}
\opt{C}{\qfigi{p1_chem_s1_q12_3}}
\opt{D}{\qfigi{p1_chem_s1_q12_4}}
\end{qoptions}

\qn{13} If the binding energy of an electron in 3$^{\text{rd}}$ excited state of a H-like atom is 2eV, then 1$^{\text{st}}$ excitation energy for that atom :
\begin{qoptions}
\opt{A}{36 eV}
\opt{B}{12 eV}
\opt{C}{48 eV}
\opt{D}{24 eV}
\end{qoptions}

\qn{14} \textbf{Statement-1 :} The hydrolysis of ethyl acetate is a second order reaction, in reality.\\
CH$_3$COOC$_2$H$_5$ + H$_2$O $\to$ CH$_3$COOH + C$_2$H$_5$OH\\
$r=K[\mathrm{CH_3COOC_2H_5}]\,[\mathrm{H_2O}]$\\
\textbf{Statement-2 :} As the concentration of water does not get altered much during the course of reaction, the reaction behaves as first order reaction with respect to ester.
\begin{qoptions}
\opt{A}{Statement-1 and Statement 2 are correct and Statement 2 explains Statement 1}
\opt{B}{Statement-1 and Statement-2 are correct and Statement -2 doesn't explain Statement 1}
\opt{C}{Statement-1 is correct and Statement-2 is incorrect}
\opt{D}{Statement-1 and Statement-2 both are incorrect.}
\end{qoptions}

\qn{15} A solution at 25$^\text{o}$C is 0.01 M in acetic acid and 0.05 M in sodium acetate. Then the pOH of solution is.\\
$\big(K_a(\text{acetic acid})=1.8\times10^{-5};$\\ $\log5=0.7;\ \log2=0.3;\ \log3=0.47\ \&\ \log3.47=0.54\big)$
\begin{qoptions}
\opt{A}{4.45}
\opt{B}{5.45}
\opt{C}{8.55}
\opt{D}{9.55}
\end{qoptions}

\qn{16} For oxidation of iron : 4Fe(s) + 3O$_2$(g) $\to$ 2Fe$_2$O$_3$(s); $\Delta_rH^\circ_{300}=-1650$ kJ mole$^{-1}$ and $\Delta_rS^\circ_{300}=-550$ JK$^{-1}$ mol$^{-1}$. Calculate $\Delta_rS^\circ_{\text{total}}$ at 300 K :
\begin{qoptions}
\opt{A}{0}
\opt{B}{+4950 J K$^{-1}$ mol$^{-1}$}
\opt{C}{$-6050$ J K$^{-1}$ mol$^{-1}$}
\opt{D}{$-1485$ J K$^{-1}$ mol$^{-1}$}
\end{qoptions}

\qn{17} 3 mol of gaseous phosphorus atoms are combined to formed P$_4$(g). Calculate $\Delta$H for this process.\\
Given B.E.(P$-$P) = 300 kJ/mol
\begin{qoptions}
\opt{A}{+1350 kJ/mol.}
\opt{B}{$-$1250 kJ/mol}
\opt{C}{$-$1050 kJ/mol}
\opt{D}{$-$1350 kJ/mol}
\end{qoptions}

\qn{18} Which of the following do not decolorize bromine water?
\begin{qoptions}
\opt{A}{2-Butene}
\opt{B}{Propane}
\opt{C}{Ethene}
\opt{D}{Cyclohexene}
\end{qoptions}

\qn{19} Mac Arthur forest process is applied for extraction of which of the following metals.\\
(I) Zn \quad (II) Au \quad (III) Cu \quad (IV) Ag \quad (V) Al
\begin{qoptions}
\opt{A}{I only}
\opt{B}{II only}
\opt{C}{IV and V only}
\opt{D}{II and IV only}
\end{qoptions}

\qn{20} For which of the element, the highest possible oxidation state shown by the element is not equal to its valence shell electron.
\begin{qoptions}
\opt{A}{Ru}
\opt{B}{Fe}
\opt{C}{V}
\opt{D}{Ti}
\end{qoptions}

\sect{Section-II}

\qn{1} 300g mixture of C$_2$H$_6$ and HCHO is burned completely with oxygen and required 400gm O$_2$. Find the mass percentage of HCHO in original sample?

\qn{2} Number of moles of grignard reagent used for 1 mole of given compound is.
\qfig{p1_chem_s2_q02_1}

\qn{3} A first order reaction completes 75\% in 40 minutes. In what time (in hours), the concentration of reactant will reduce from 0.64 M to 0.01 M ?

\qn{4} The reactant `A' undergoes intramolecular aldol condensation reaction when heated in presence of NaOH to obtain the product `B'
\qfig{p1_chem_s2_q04_1}
The number of methylene ($-$CH$_2-$) units in `A' is

\qn{5} The total possible number of co-ordination isomers possible for $[\mathrm{Mo(en)_2Br_2}][\mathrm{AgF_4}]$

\partsec{3}{MATHEMATICS}
\sect{Section-I}

\qn{1} If the equation $\left|x^2+4x+3\right|-mx+2m=0$ has exactly three solutions, then the value(s) of $m$ is equal to :-
\begin{qoptions}
\opt{A}{$-8\pm2\sqrt{15}$}
\opt{B}{$-8-2\sqrt{15}$}
\opt{C}{3}
\opt{D}{$-8+2\sqrt{15}$}
\end{qoptions}

\qn{2} The number of solution(s) of the equation $z^2=4z+|z|^2+\dfrac{16}{|z|^3}$ is\\
(where $z=x+iy$, $x,y\in R$, $i^2=-1$ and $x\neq2$)
\begin{qoptions}
\opt{A}{0}
\opt{B}{1}
\opt{C}{2}
\opt{D}{3}
\end{qoptions}

\qn{3} \textbf{Assertion (A) :} Slope of the curve given as $y=x^2$ at $x=1$ is not defined.\\
\textbf{Reason (R) :} Slope of the curve given as $y=x^2$ at $x=1$ is $\pm\dfrac12$.
\begin{qoptions}
\opt{A}{Both A and R are true and R is the correct explanation of A}
\opt{B}{Both A and R are true but R is NOT the correct explanation of A}
\opt{C}{A is true but R is false.}
\opt{D}{Both A and R are false.}
\end{qoptions}

\qn{4} Consider two sets A $=\{1,2,3,4,5\}$ and B $=\{1,2,3,\dots..9,10\}$. Out of all the functions which can be defined from A $\to$ B, one functions randomly selected. The probability that the functions selected is injective, lies in the interval.
\begin{qoptions}
\opt{A}{$\left(\dfrac3{25},\dfrac4{25}\right)$}
\opt{B}{$\left(\dfrac4{25},\dfrac5{25}\right)$}
\opt{C}{$\left(\dfrac5{25},\dfrac6{25}\right)$}
\opt{D}{$\left(\dfrac7{25},\dfrac8{25}\right)$}
\end{qoptions}

\qn{5} If $\cos\theta=\dfrac{ab+bc+ca}{a^2+b^2+c^2}$, where $a,b,c$ are sides of a $\Delta$ABC, then $\theta$ can not be equal to :-
\begin{qoptions}
\opt{A}{$\dfrac\pi{12}$}
\opt{B}{$\dfrac\pi6$}
\opt{C}{$\dfrac\pi4$}
\opt{D}{$\dfrac{5\pi}{12}$}
\end{qoptions}

\qn{6} Let $S_n=\displaystyle\sum_{k=0}^n{}^{n+k}C_k\,k$.\\
If $S_{100}=\dfrac{A!}{B!}\left(\dfrac1{101!}-\dfrac1{102!}\right)$ then the value of $A+B$ is
\begin{qoptions}
\opt{A}{150}
\opt{B}{300}
\opt{C}{450}
\opt{D}{600}
\end{qoptions}

\qn{7} There are two groups of students A and B in a sports academy which offers three games volleyball, basketball, baseball. The ratio of students playing volleyball, basketball and baseball is 3:4:5 in both the group. If the ratio of total number of students in Group A and B is 2:3 and there are 40 more students in Group B than Group A who plays baseball, then the total number of students in the sports academy is
\begin{qoptions}
\opt{A}{360}
\opt{B}{720}
\opt{C}{480}
\opt{D}{960}
\end{qoptions}

\qn{8} A boat rowed with velocity $V$ directly across a stream of width `$b$'. If velocity of current is directly proportional to product of distances from two banks, find the distance downstream to the point where it lands.
\begin{qoptions}
\opt{A}{$\dfrac{kb^2}{2V}$}
\opt{B}{$\dfrac{kb^3}{2V}$}
\opt{C}{$\dfrac{kb^3}{6V}$}
\opt{D}{$\dfrac{kb^2}{6V}$}
\end{qoptions}

\qn{9}
\par\nobreak
\begin{center}\small
\begin{tabular}{|c|p{4.2cm}|c|c|}\hline
 & \multicolumn{1}{c|}{\textbf{List-I}} & & \textbf{List-II}\\\hline
(A) & If $x\dfrac{dy}{dx}=x^2+y-2$, $y(1)=1$, then $y(2)$ equals & (I) & 1\\\hline
(B) & Let $S=$\newline\resizebox{\linewidth}{!}{$\displaystyle\sum_{n=1}^{9999}\frac1{\left(\sqrt n+\sqrt{n+1}\right)\left(\sqrt[4]n+\sqrt[4]{n+1}\right)}$}\newline then $\sqrt S$ equals & (II) & 2\\\hline
(C) & Let $S$ be the area bounded by the curve $y=\sin x$ ($0\leq x\leq\pi$) and the $x$-axis and $T$ be the area bounded by the curves $y=\sin x\left(0\leqslant x\leqslant\dfrac\pi2\right)$, $y=a\cos x\left(0\leqslant x\leqslant\dfrac\pi2\right)$ and the $x$-axis (where $a\in R^+$). The value of $(3a)$ such that $S:T=1:\dfrac13$, is & (III) & 3\\\hline
(D) & If the curve C in the $xy$ plane has the equation $x^2+xy+y^2=1$, then the square of greatest distance of a point on the curve C from the origin, is & (IV) & 4\\\hline
\end{tabular}
\end{center}
\begin{qoptions}
\opt{A}{(A)-II; (B)- III; (C)-IV ; (D)-I}
\opt{B}{(A)-I; (B)-III; (C)-IV ; (D)-II}
\opt{C}{(A)-II; (B)-III; (C)-IV ; (D)-II}
\opt{D}{(A)-III; (B)-II; (C)-IV ; (D)-I}
\end{qoptions}

\qn{10} Sum of digits of the number of rational terms in the expansion of $\left(\sqrt2+\sqrt3+\sqrt[3]5\right)^{20}$, is :-
\begin{qoptions}
\opt{A}{9}
\opt{B}{8}
\opt{C}{7}
\opt{D}{6}
\end{qoptions}

\qn{11} For an ellipse $\dfrac{x^2}9+\dfrac{y^2}4=1$, with vertices A \& A$'$, tangent drawn at the point `P' in the first quadrant meets the $y$-axis in Q and the chord A$'$P meets the $y$-axis in M. If `O' is the centre of ellipse, then OQ$^2-$MQ$^2$ is equal to,
\begin{qoptions}
\opt{A}{9}
\opt{B}{13}
\opt{C}{4}
\opt{D}{5}
\end{qoptions}

\qn{12} Let A $=\begin{bmatrix}3x^2\\1\\6x\end{bmatrix}$, B $=[a\ b\ c]$ and C $=\begin{bmatrix}(x+2)^2&5x^2&2x\\5x^2&2x&(x+2)^2\\2x&(x+2)^2&5x^2\end{bmatrix}$ be three given matrices, where $a,b,c$ and $x\in R$, Given that tr.(AB) $=$ tr.(C) $\forall\ x\in R$, (where tr.(X) denotes trace of square matrix X).\\
If $\displaystyle\int_0^\infty\dfrac{\ln x}{cx^2+ax+b}dx=\dfrac{\pi\ln2}{\sqrt q}$, then the value of $q$ is
\begin{qoptions}
\opt{A}{29}
\opt{B}{27}
\opt{C}{26}
\opt{D}{28}
\end{qoptions}

\qn{13} Find the value of the definite integral\\
$2^{2010}\dfrac{\displaystyle\int_0^1x^{1004}(1-x)^{1004}dx}{\displaystyle\int_0^1x^{1004}\left(1-x^{2010}\right)^{1004}dx}$.
\begin{qoptions}
\opt{A}{1005}
\opt{B}{2010}
\opt{C}{4020}
\opt{D}{8040}
\end{qoptions}

\qn{14} A line L passing through the point P (1, 4, 3), is perpendicular to both the lines $\dfrac{x-1}2=\dfrac{y+3}1=\dfrac{z-2}4$ and $\dfrac{x+2}3=\dfrac{y-4}2=\dfrac{z+1}{-2}$. If the position vector of point Q on L is $(a_1,a_2,a_3)$ such that $(PQ)^2=357$, then $(a_1+a_2+a_3)$ can be:
\begin{qoptions}
\opt{A}{2}
\opt{B}{3}
\opt{C}{15}
\opt{D}{16}
\end{qoptions}

\qn{15} $f(x)=\left|\{x\}-\dfrac12\right|$ \& $g(x)=\sin x$ where $\{.\}$ is fractional part.\\
\textbf{Statement I:} The composite function $f(g(x))$ is non-differentiable in $x\in(0,2\pi)$ at 5 points.\\
\textbf{Statement II:} If $f(x)$ is non-differentiable \& $g(x)$ is differentiable, then $f(g(x))$ is non-differentiable function.
\begin{qoptions}
\opt{A}{Both statements are true, statement-II is correct explanation of statement-I.}
\opt{B}{Both statement are true, statement-II is not correct explanation of statement-I}
\opt{C}{First statement true \& II statement false.}
\opt{D}{First statement false and II statement true.}
\end{qoptions}

\qn{16} If $\log^2y+\left(2^{1+x}+2^{1-x}\right)(\log y)+2^{2x+1}+2^{1-2x}=0$; where $x,y\in R$, then the value of $\log(x+y)$ is :- (where base of log is e)
\begin{qoptions}
\opt{A}{$-1$}
\opt{B}{1}
\opt{C}{$-2$}
\opt{D}{2}
\end{qoptions}

\qn{17} A line through the origin meets the circle $x^2+y^2=a^2$ at P and the hyperbola $x^2-y^2=a^2$ at Q. Then locus of the point of intersections of tangent to the circle at P with the tangent at Q to the hyperbola is :-
\begin{qoptions}
\opt{A}{$\left(a^6+2y^4\right)x^2=a^4$}
\opt{B}{$\left(a^4+2y^4\right)x^2=a^3$}
\opt{C}{$\left(a^6+4y^4\right)x^2=a^4$}
\opt{D}{$\left(a^4+4y^4\right)x^2=a^6$}
\end{qoptions}

\qn{18} Area of the region in $x$-$y$ plane satisfying $x^{10}-x^4+y^2\leq0$, is $A$, then $[A]$ (where $[x]$ represents greatest integer function)
\begin{qoptions}
\opt{A}{1}
\opt{B}{2}
\opt{C}{3}
\opt{D}{4}
\end{qoptions}

\qn{19} Given $\vec a,\vec b,\vec c$ are vectors such that $\left[\vec a\ \vec b\ \vec c\right]=\dfrac13$. If the vector $\vec V$ can be expressed as linear combination of $\vec b\times\vec c$, $\vec c\times\vec a$ and $\vec a\times\vec b$ as\\
$\vec V=x(\vec b\times\vec c)+y(\vec c\times\vec a)+z(\vec a\times\vec b)$ then $(x+y+z)$ has the value equal to
\begin{qoptions}
\opt{A}{$\vec V\cdot\left(\vec a+\vec b+\vec c\right)$}
\opt{B}{$3\vec V\cdot\left(\vec a+\vec b+\vec c\right)$}
\opt{C}{$2\vec V\cdot\left(\vec a+\vec b+\vec c\right)$}
\opt{D}{None}
\end{qoptions}

\qn{20} The value of $\displaystyle\lim_{n\to\infty}\sum_{r=1}^n\dfrac{1.3.5.\dots\dots.(2r-1)}{2.4.6.\dots\dots.(2r+2)}$, is (where $n\in N$)
\begin{qoptions}
\opt{A}{$\dfrac12$}
\opt{B}{0}
\opt{C}{1}
\opt{D}{Does not exist}
\end{qoptions}

\sect{Section-II}

\qn{1} Two friends have equal number of sons. There are 3 tickets for a cricket match which are to be distributed among the sons. The probability that two tickets go to the sons of one and one tickets goes to the sons of other is $\dfrac67$. Then total number of boys is equal to (sum of son of each friend).

\qn{2} Let $f(x)=\cos x\left(\sin x+\sqrt{\sin^2x+\sin^2\theta}\right)$ where $\theta$ is given constant, then maximum value of $f(x)$ is $g(\theta)$. The maximum value of $g^2(\theta)$, is

\qn{3} Let $f(x)$ be a non-constant thrice differentiable function defined on $(-\infty,\infty)$ such that $f(x)=f(6-x)$ and $f'(0)=0=f'(2)=f'(5)$. Determine the minimum number of zeroes of $g(x)=\left(f''(x)\right)^2+f'(x)f'''(x)$ in the interval $[0,6]$.

\qn{4} Let A $=\sin\dfrac\pi{81}.\sin\dfrac{2\pi}{81}\dots\sin\dfrac{80\pi}{81}$ and\\
B $=\sin\dfrac\pi{27}.\sin\dfrac{2\pi}{27}\dots\sin\dfrac{26\pi}{27}$. If the value of the following expression $\sqrt[7]{\dfrac{12B}A}$ is $2^k$, then $k$ is

\qn{5} Given $2x-y+2z=1$; $x-2y+z=-4$; \& $x+y+\lambda z=4$. Then the value of $\lambda$ such that the given system of equation has no solution is

\end{multicols}
\end{document}
